\section{Products with smooth backgrounds}
\label{sec:psh-bounded-products}

The Bedford--Taylor construction also defines products after adding smooth forms to complex Hessians. Independence of the smooth correction makes these products global; positivity requires retaining the original lower curvature bound. Fix $1\leq k\leq n$.

\Needspace{8\baselineskip}
Let $u$ be locally bounded and $\theta$-plurisubharmonic. On a relatively compact coordinate ball with coordinates $z$, choose $A>0$ and put $q=A|z|^2$ so that $p=u+q$ is plurisubharmonic. Set $\beta=\theta-\ddc q$ and define
\begin{equation}\label{eq:bounded-expansion}
 (\theta+\ddc u)^k
 :=\sum_{j=0}^k\binom{k}{j}(\ddc p)^j\wedge\beta^{k-j}.
\end{equation}
Each summand is defined by \cref{thm:bt-input} and multiplication by a smooth form. The signed terms in this expansion should not individually be assumed positive.

\begin{proposition}\label{prop:bounded-products}
The construction \eqref{eq:bounded-expansion}, and its mixed version, is independent of the local correction $q$ and of the decomposition of each current into a smooth form and a $\ddc$-potential. It defines global products of currents with locally bounded quasi-plurisubharmonic potentials. The products are symmetric, multilinear, plurifinely local, and give no mass to pluripolar sets. If every factor is positive, their mixed product is positive.
\end{proposition}
\begin{proof}
On the overlap of two coordinate balls, the two corrected potentials differ by a smooth function. The multilinear rule
$\ddc(p+f)=\ddc p+\ddc f$, valid for bounded local products by smooth approximation, and the binomial formula identify the two expansions. The same argument handles a smooth change of the decomposition of a current. Any two potential decompositions of the same $(1,1)$-current do differ smoothly: their difference has smooth $\ddc$, and taking its local Euclidean trace gives smoothness by elliptic regularity. Symmetry, multilinearity and locality then follow term by term. Wedges of bounded Bedford--Taylor products with smooth forms have locally finite variation and do not charge pluripolar sets, so these properties also pass to the finite expansion.

It remains to prove positivity; it does not follow from the signs of the summands. Write the factors as $T_i=\theta_i+\ddc u_i$, $1\leq i\leq k$, where each $\theta_i$ is smooth and real and each $u_i$ is locally bounded and quasi-plurisubharmonic. Assume $T_i\geq0$. Fix $\varepsilon>0$ and a local positive constant form $\kappa$. After shrinking a neighborhood of a chosen point, each smooth background $\theta_i$ admits a constant, hence closed, form $\theta_i'$ such that
\[
 -\varepsilon\kappa\leq\theta_i-\theta_i'
 \leq\varepsilon\kappa.
\]
Then
$S_i=\theta_i'+\varepsilon\kappa+\ddc u_i$ is positive and closed there, with a bounded local plurisubharmonic potential. Moreover
\[
 T_i+2\varepsilon\kappa
 =S_i+(\theta_i-\theta_i'+\varepsilon\kappa).
\]
All terms on the right are positive, and the second term is smooth. Expanding the mixed product therefore gives a sum of positive currents. Although the neighborhoods depend on $\varepsilon$, they cover the original coordinate ball. The product of the $T_i+2\varepsilon\kappa$ is thus positive throughout that ball. By multilinearity it converges weakly to the product of the $T_i$ as $\varepsilon\to0$, proving positivity.
\end{proof}

The local expansion also gives continuity along uniformly locally bounded decreasing sequences, by applying \cref{thm:bt-input} to the corrected potentials and then summing the finitely many terms. Each bounded top-degree product has finite mass on compact $X$. For bounded $u\in\PSH(X,\theta)$, we write $\mathrm{MA}_\theta(u)=(\theta+\ddc u)^n$.

For $C^{1,1}$ potentials, these measure-valued products agree with the products of their measurable Hessian coefficients.

\begin{proposition}\label{prop:c11-products}
Let $u_1,\ldots,u_k$ be locally $C^{1,1}$ quasi-plurisubharmonic functions and let $\theta_1,\ldots,\theta_k$ be smooth real $(1,1)$-forms. The bounded product
\[
 (\theta_1+\ddc u_1)\wedge\cdots\wedge(\theta_k+\ddc u_k)
\]
coincides with the almost-everywhere wedge of the locally bounded measurable coefficients of its factors.
\end{proposition}
\begin{proof}
On a smaller coordinate ball, add a common smooth quadratic $q$ so that each $u_i+q$ is plurisubharmonic, and convolve with a smooth radial probability kernel. The Hessians of the convolutions are uniformly locally bounded and converge to the weak Hessians in every finite $L^s$ space. Their finite products therefore converge in $L^1$ on smaller balls. The potentials converge uniformly, so bounded Bedford--Taylor continuity identifies this limit with their pluripotential product. Expanding off the smooth correction, as in \eqref{eq:bounded-expansion}, proves the assertion for the original factors.
\end{proof}

For a locally bounded quasi-plurisubharmonic function $v$ and a smooth convex increasing function $F$ on a neighborhood of its range, the bounded Bedford--Taylor chain rule is
\begin{equation}\label{eq:bounded-convex-chain-rule}
 \ddc F(v)=F'(v)\ddc v+F''(v)\,\mathrm{d}v\wedge\dc v.
\end{equation}
The gradient term and its products with positive bounded Bedford--Taylor factors are nonnegative. These local identities follow by decreasing regularization; bounded potentials have the local Sobolev regularity needed for the gradient expression. They are part of the bounded local calculus recalled in \cite{SGPT,GZ17}. In particular, inequalities obtained from this formula pass through local smoothing with a lower-curvature error tending to zero.
